% =====================================================================
\section{Audio Processing Pipeline}
\label{sec:audio}
% =====================================================================

The audio pipeline turns a waveform into a sequence of acoustic feature
vectors at the video frame rate. Unlike the visual crops, audio features are
computed on the fly during training, because noise is mixed into each
utterance afresh every time it is loaded (\cref{sec:noise}).

\subsection{From waveform to features}
\paragraph{Loading and resampling.} Waveforms are read with \texttt{soundfile}
as 32-bit floating-point arrays in $[-1,1]$. Stereo input is averaged to mono,
and any input that is not at 16\,kHz is resampled with
\texttt{torchaudio.functional.resample}. GRID audio is resampled once during
preprocessing; microphone audio is captured at 16\,kHz directly where the
device allows it.

\paragraph{Amplitude normalisation.} The waveform presented to the model---the
noisy mixture during training, the raw microphone signal at inference---is
scaled to a fixed RMS level of $-25$\,dBFS \rec, with the gain limited to
$+20$\,dB so that near-silence is not amplified into loud noise. This removes
differences in recording level between GRID speakers and the demonstration
microphone. Normalisation is applied \emph{after} noise mixing so that it
does not alter the SNR.

\paragraph{Framing and windowing.} The signal is split into 25\,ms frames
(400 samples) with a 10\,ms hop (160 samples), and each frame is multiplied by a
Hann window to reduce spectral leakage \rec. These values are the conventional
choice for speech, because speech is approximately stationary over 20--30\,ms.
The 10\,ms hop yields 100 feature frames per second, exactly four per video
frame.

\paragraph{Log-Mel spectrogram.} A 512-point short-time Fourier transform
(STFT) gives the power spectrum $|X_k(\tau)|^2$ of each frame $\tau$. A bank of
$M=80$ triangular Mel filters $H_m(k)$ spanning 20\,Hz--8\,kHz \rec\ is applied,
and the result is log-compressed:
\begin{equation}
\mathrm{LM}_m(\tau) = \log\!\Big(\epsilon + \sum_{k} H_m(k)\,|X_k(\tau)|^2\Big),
\qquad m = 1,\dots,80,\quad \epsilon = 10^{-6}.
\label{eq:logmel}
\end{equation}
The Mel scale spaces filters more densely at low frequencies, mirroring human
pitch perception, and the logarithm compresses the large dynamic range of
speech energy.

\paragraph{MFCCs.} Mel-frequency cepstral coefficients are obtained by applying a
discrete cosine transform (DCT) to the log-Mel vector and keeping the first 13
coefficients, optionally with their first and second temporal derivatives
($\Delta$, $\Delta\Delta$) to give 39 values per frame:
\begin{equation}
c_j(\tau) = \sum_{m=1}^{M} \mathrm{LM}_m(\tau)\cos\!\Big[\frac{\pi j}{M}\Big(m-\tfrac12\Big)\Big],
\qquad j = 0,\dots,12.
\end{equation}

\paragraph{Feature normalisation.} Each feature dimension is standardised with a
mean and standard deviation estimated on the training partition, under the same
SNR distribution used in training. The statistics are computed once by
\texttt{compute\_stats.py}, stored in \texttt{stats/feature\_norm.npz} and
loaded unchanged at validation, test and inference time. Statistics are never
computed from validation or test data.

\subsection{Choice of audio features}
The proposal lists both log-Mel spectrograms and MFCCs \spec\ but does not
state whether they are used together. Because MFCCs are a linear transform
(the DCT) of the log-Mel features, a neural encoder given log-Mel input can in
principle learn any information that MFCCs contain. The plan therefore uses a
manageable initial configuration and treats the others as experiments:

\begin{table}[H]
\centering
\caption{Audio feature configurations.}
\label{tab:audio-configs}
\small
\begin{tabularx}{\textwidth}{@{}llX@{}}
\toprule
\textbf{ID} & \textbf{Features (dimension)} & \textbf{Role} \\
\midrule
A1 & 80 log-Mel & Primary configuration for all models \rec \\
A2 & 13 MFCC + $\Delta$ + $\Delta\Delta$ (39) & Experiment: compact, decorrelated input \\
A3 & log-Mel $\oplus$ MFCC (119) & Experiment: the combined input named in the proposal \\
\bottomrule
\end{tabularx}
\end{table}

All three configurations are computed by the same \texttt{torchaudio} transform
objects, so switching between them is a one-line configuration change. A2 and
A3 are run on the audio-only baseline first, because it trains fastest; the
best configuration is then adopted for the fused model if it improves
validation WER.

\subsection{Audio encoder}
\label{sec:audio-encoder}
The audio encoder maps the 100\,fps features to a 25\,fps sequence of acoustic
representations aligned with the visual sequence. Two 1-D convolutions with
stride~2 reduce the frame rate by a factor of four while mixing information
across neighbouring frames. A two-layer BiLSTM, the encoder specified in the
proposal \spec, then models context across the whole utterance. Strided
convolutions are preferred to simply discarding frames, because they learn
which local acoustic detail to retain.

For a convolution with kernel $k=3$, padding $p=1$ and stride $s=2$, an input of
length $L$ produces an output of length $\lfloor (L + 2p - k)/s \rfloor + 1 =
\lceil L/2 \rceil$. Two such layers turn $T_a$ audio frames into $\lceil
T_a/4\rceil$ encoder frames. The same formula is used in the data loader to
compute per-utterance output lengths for CTC (\cref{sec:ctc}).

\subsection{Tensor-shape walkthrough}
\Cref{tab:audio-shapes} assumes a batch of $B=16$ clips of 3\,s at 16\,kHz
($N=48{,}000$ samples) and configuration A1. The STFT with centred frames
produces $N/160+1=301$ frames; the last frame is dropped so that $T_a = 300 =
4T$ \rec.

\begin{table}[H]
\centering
\caption{Audio encoder tensor shapes for an example batch.}
\label{tab:audio-shapes}
\small
\begin{tabularx}{\textwidth}{@{}lXl@{}}
\toprule
\textbf{Stage} & \textbf{Operation} & \textbf{Output shape} \\
\midrule
Waveform & Normalised 16\,kHz mono audio & \shape{16,48000} \\
Features & Log-Mel (\cref{eq:logmel}), trim to $4T$, standardise & \shape{16,300,80} \\
Transpose & Channels-first layout for Conv1d & \shape{16,80,300} \\
Subsample 1 & Conv1d($80{\to}256$, $k{=}3$, $s{=}2$), BN, ReLU & \shape{16,256,150} \\
Subsample 2 & Conv1d($256{\to}256$, $k{=}3$, $s{=}2$), BN, ReLU & \shape{16,256,75} \\
Transpose & Time-first layout for LSTM & \shape{16,75,256} \\
Temporal & 2-layer BiLSTM, 256 units per direction, dropout 0.3 & \shape{16,75,512} \\
\bottomrule
\end{tabularx}
\end{table}

The dimensions are: batch ($B$), time (audio frames at 100\,fps before
subsampling, encoder frames at 25\,fps after), and features (Mel bands, then
convolution channels, then concatenated forward and backward LSTM states). The
output has the same shape as the visual encoder output, which simplifies
fusion.

\subsection{Variable-length sequences and batching}
Most GRID utterances have the same 3\,s duration, but some are shorter, and
windows in the live demo may vary. The collate function therefore
\begin{enumerate}
  \item pads every waveform in a batch with zeros to the longest one,
  \item records the true number of feature frames $T_a^{(b)}$ for each item $b$,
  \item computes encoder lengths $T^{(b)} = \lceil T_a^{(b)}/4 \rceil$, and
  \item builds a Boolean padding mask of shape \shape{B,T}.
\end{enumerate}
The BiLSTM receives the lengths through \texttt{pack\_padded\_sequence}, so padded
frames do not contaminate the backward direction; the fusion attention ignores
padded keys through the mask; and CTC receives the true lengths. Utterances are
grouped into batches of similar length (bucketing) when lengths vary widely, to
reduce wasted computation on padding.

\subsection{Testing the audio pipeline}
\begin{itemize}
  \item The log-Mel output of \texttt{torchaudio} is compared with
        \texttt{librosa.feature.melspectrogram} on a test file using identical
        parameters; the maximum absolute difference in log-energy must be below a
        small tolerance (for example $10^{-3}$).
  \item A 3\,s waveform must produce exactly 300 feature frames and 75 encoder
        frames.
  \item Normalised training features must have per-band mean within 0.05 of 0
        and standard deviation within 0.05 of 1.
  \item The audio encoder on a random input of \shape{2,300,80} returns
        \shape{2,75,512}.
\end{itemize}
